% ============================================================
%  导读
% ============================================================
\clearpage
\section*{导读：一条走了四十七年的暗线}
\addcontentsline{toc}{section}{导读：一条走了四十七年的暗线}
\markboth{导读}{导读}

\chapterquote{“Atiyah explained to me that what Goddard, Nuyts, and Olive were defining as the ‘magnetic gauge group’ was actually the same as the Langlands dual group in the context of the Langlands program of number theory. I am pretty sure that Atiyah is the one who made this connection, though as far as I know he never wrote anything about it.”}{Edward Witten, 2021}

\subsection*{一、1977 年春天}

1977 年春天，Michael Atiyah 在哈佛访问，讲的是瞬子——四维 Yang--Mills 方程的自对偶解。听讲的人里有一位 25 岁的 Junior Fellow，Edward Witten，此前做的是粒子物理与广义相对论。Atiyah 的讲法显然打动了他：讲座之后两人开始交谈，Witten 在当年年底受邀回访牛津。那一个月里，Atiyah 正在讲四维瞬子——他与 Vladimir Drinfeld、Nigel Hitchin、Yuri Manin 刚刚借助扭量理论与代数几何发现了 ADHM 构造。瞬子当时是物理里最热的话题之一，很大程度上因为 Gerard 't Hooft 用它解决了核力的“$U(1)$ 问题”；到 1977 年底，人们又从 Albert Schwarz 的观察中知道，'t Hooft 的这项工作可以理解为 Atiyah--Singer 指标定理的一次应用。

到了那个月的月底，Atiyah 忽然转到一个完全不同的题目。他拿出两篇 Witten 没见过的论文：Goddard--Nuyts--Olive 的《规范理论与磁荷》，和 Montonen--Olive 的《磁单极子作为规范粒子？》。他说这里面有很深的东西，Witten 应当去伦敦找 Olive 谈谈。

Witten 当时没听说过 Langlands 纲领，也不知道 Robert Langlands 是谁。但 Atiyah 说，这个数论纲领非常深刻，物理学家必然也在通向某种重要的东西。四十四年后，Witten 在 Isaac Newton Institute 的这次演讲，第一次把这条线索系统地讲成一个故事。

\subsection*{二、这场演讲在讲什么}

题目里有三个名词：对偶、Langlands 对应、Atiyah 的影响。对应地，演讲有三条线索，而它们最终拧成一股。

\textbf{第一条是电磁对偶。}1977 年的 GNO 论文指出：把规范群 $G$ 希格斯化到极大环面 $T$ 之后，电电荷由 $T$ 的特征标（即表示）分类，磁电荷由 $T$ 的余特征标分类；而 $G$ 与它的对偶群 $\gvee$ 之间，恰好交换这两件事。Montonen 与 Olive 更进一步：他们猜测，两个规范群互为对偶的量子规范理论根本是同一个理论，只是强耦合与弱耦合互换。这个猜想在 1977 年近乎狂想；要让它站住，需要超对称、需要 Bogomolny 方程、需要 Atiyah--Hitchin 的单极子模空间、需要 Sen 的 $L^2$ 调和形式，最后在 Seiberg--Witten 的工作里变成理解强耦合动力学的钥匙。

\textbf{第二条是几何 Langlands。}1990 年前后，Beilinson 与 Drinfeld 把 Langlands 纲领从数域搬到 Riemann 面上，是为几何 Langlands。它的核心对象是 Hecke 变换：在一点修改一个全纯丛。Witten 多年读不懂这套语言，尤其不懂 Hecke 变换怎么会和物理有关。答案是：Hecke 变换就是量子规范论里的 't Hooft 算子，而它的“类型”由对偶群的表示标号——这正是 1977 年 GNO 那句话的另一种说法。

\textbf{第三条是 Atiyah 的方法。}整场演讲里，Atiyah 本人只出现在开头与结尾，但每一个转折处都有他的影子：把单极子看成非线性椭圆方程的解，把 Donaldson 理论看成扭曲的超对称规范论，把四维流形的拓扑问题交给物理去算，以及那句“这个联系我从没写下来过”。Witten 没有明说，但读者不难看出：这场演讲名义上在讲对偶，实际上在讲一种做数学与物理的方法——先相信两个互不相干的领域在描述同一个结构，再去把它找出来。

\subsection*{三、为什么值得读}

对物理读者，这是一个关于“猜想的生命史”的案例：Montonen--Olive 猜想从 1977 年被提出，到 1979 年被 Osborn 修正，到 1994 年被 Sen 用一个可计算的检验钉住，再到同年成为 Seiberg--Witten 理论的出发点。它展示了理论物理里最罕见的一类进步：不是实验，而是一个内在一致性的要求逼出了一个新结构。

对数学读者，这是一个关于“同一性”的案例：根数据的对偶（纯代数）、磁荷与电电荷的交换（量子场论）、Hecke 特征层与局部系统的对应（几何表示论），三件事在各自的领域里被独立发现，最后被证明是同一件事的不同侧面。

对一般读者，这是一份极好的历史入口：54 页幻灯片里出现了二十多个名字、十几篇论文，每一篇都值得单独读；本书的批注与附录，就是为这些岔路准备的地图。

\subsection*{四、年表}

\begin{center}\small
\begin{longtable}{@{}p{2.1cm}p{10.9cm}@{}}
\toprule
\textbf{年份} & \textbf{事件} \\
\midrule
\endhead
1969 & Narasimhan 与 Ramanan 在紧 Riemann 面上构造 Hecke 算子，几何 Langlands 的技术前身。\\
1974 & 't Hooft 与 Polyakov 各自指出：非阿贝尔规范理论希格斯化后，磁单极子是光滑的孤子解。\\
1975 & Bogomolny 给出能量下界（BPS 界）；Prasad 与 Sommerfield 写出第一个显式解。\\
1977.1 & Goddard、Nuyts、Olive 的《规范理论与磁荷》发表，提出“磁规范群”。\\
1977 春 & Atiyah 在哈佛把 GNO 的磁规范群指认为 Langlands 对偶群。\\
1977 & Montonen 与 Olive 提出电—磁对偶：规范群为 $G$ 与 $\gvee$ 的两个理论量子等价。\\
1978 & Witten 与 Olive 发现超对称代数中的中心荷，证明 BPS 质量公式在量子上不被修正。\\
1979 & Osborn 指出只有 $\cN=4$ 超对称理论才具备 Montonen--Olive 对偶所需的性质。\\
1988 & Atiyah 与 Hitchin 的磁单极子专著出版；Witten 指出 Donaldson 理论是扭曲的 $\cN=2$ 理论。\\
1990 前后 & Beilinson 与 Drinfeld 提出几何 Langlands 纲领。\\
1994.2 & Sen 用两单极子模空间上的 $L^2$ 调和形式检验对偶，给出“确凿证据”。\\
1994.7 & Seiberg 与 Witten 解出 $\cN=2$ 超对称 Yang--Mills 的低能有效理论。\\
1994.10 & Witten 把 $\cN=2$ 理论扭曲，得到 Seiberg--Witten 方程，重证 Donaldson 不变量。\\
1990 年代 & Lusztig、Ginzburg、Mirković、Vilonen 等建立几何 Satake 等价。\\
2005 & Kapustin 给 't Hooft 算子一个更精确的奇点定义。\\
2006 & Kapustin 与 Witten 把几何 Langlands 解释为四维超对称规范理论的 S 对偶。\\
2009 & Witten 从六维 $(2,0)$ 理论重新导出这一图像。\\
2021.9 & 本演讲。\\
2024 & Gaitsgory 与十余位合作者发布五篇论文，证明几何 Langlands 猜想。\\
\bottomrule
\end{longtable}
\end{center}

\subsection*{五、怎么读这本书}

本书九章对应底本 54 页幻灯片，章末的“对应底本”一栏给出页码区间。三种读法：

\begin{itemize}
  \item \textbf{想先懂物理}：读第 1--3 章（GNO 与对偶群），跳过第 5 章的模空间细节，直接读第 6 章（Sen 与 Seiberg--Witten）。
  \item \textbf{想先懂数学}：读第 7--9 章（几何 Langlands、线算子、$\RR\times C$），再回头补第 3 章的根数据对偶。
  \item \textbf{想先懂历史}：读本章的年表，再读第 1、4、6 章的叙事，最后读第 9 章的收束。
\end{itemize}

阅读本书不需要 Langlands 纲领的数论背景，也不需要掌握超对称场论的计算。所需的群论（根数据、特征标格、余特征标格）在第 3 章的批注 3.1 中从头讲起；所需的规范场论（联络、曲率、单极子）在第 2 章；超对称只在第 4、6 章作为“为什么必须是 $\cN=4$、$\cN=2$”的理由出现，读者若只想知道结论，可以只读每节开头的段落。
